% ============================================================
% Thema 5: Global Optimization (S. 47-56)
% Include-Datei fuer zusammenfassung.tex
% ============================================================
\section{Thema 5: Globale Optimierung (Global Optimization) [S.~47--56]}

\subsection{Motivation und Einordnung [S.~47]}

Kapitelmotto: \glqq And then there were many.\grqq{} --- Im vorigen Kapitel wurden
Newtons Methode und der Gradientenabstieg (gradient descent) für die Nullstellensuche
entwickelt. Diese Verfahren sind jedoch \emph{lokale} Methoden: Sie konvergieren zu
derjenigen Lösung, die in der Nähe des Startpunkts liegt. Was aber, wenn diese nahe
Lösung schlecht ist --- oder eine viel bessere Lösung weiter entfernt liegt? Das Skript
formuliert es so: \glqq Time to do away with oversimplifications.\grqq{} In diesem Kapitel
lernen wir Shekel's Foxholes als klassische Testfunktion für globale Optimierung kennen,
sehen, wie sich ein Optimierungsproblem als Nullstellenproblem (root-finding problem)
formulieren lässt, und führen Strategien zum Finden \emph{globaler} Minima ein.

Der Bezug zum maschinellen Lernen ist zentral: Im Deep Learning ist das Verständnis
multidimensionaler und globaler Optimierung entscheidend, damit gro{\ss}e Modelle effektiv
trainieren, konvergieren und tatsächlich etwas Nützliches lernen. Die Loss-Landschaften
neuronaler Netze (neural network loss landscapes, Li et al.\ 2018, Fu{\ss}note~8 [S.~47])
sind hochgradig nicht-konvex, besitzen flache Regionen und enthalten viele lokale
Minima --- verschiedene Initialisierungen führen zu verschiedenen Endlösungen.

\subsection{Lernziele (Learning Objectives) [S.~50]}
\begin{enumerate}[itemsep=1pt]
  \item Erklären können, warum lokale Methoden das Finden globaler Minima nicht
        garantieren können.
  \item Strategien zum Finden globaler Extrema anwenden können.
  \item Verstehen, wie der Gradientenabstieg (gradient descent) hilft, die Optimierung
        in Richtung von Minima zu lenken.
\end{enumerate}

\subsection{Hands On: Shekel's Foxholes [S.~48--50]}

\begin{defbox}{Shekel's Foxholes (1D-Testfunktion), Gl.~(35) [S.~48]}
Shekel's Foxholes ist eine klassische Testfunktion für globale Optimierung mit
\emph{kontrollierter Multimodalität} (controlled multimodality), d.\,h.\ man kann gezielt
einstellen, wie viele Minima die Funktion hat und wie sie aussehen. In 1D lautet die
einfache Form (Original-Gl.~(35); Fu{\ss}note~9: J.~Shekel, \glqq Test functions for
multimodal search techniques\grqq, 1971):
\[
  f(x) \;=\; -\sum_{i=1}^{m} \frac{c_i}{(x - a_i)^2 + r_i} \tag{35}
\]
Dabei sind $a_i$ die Positionen der \glqq Fuchslöcher\grqq{} (foxholes), $c_i$ steuert die
\emph{Tiefe} jedes Lochs und $r_i$ die \emph{Breite}. Durch Anpassen dieser Parameter
lässt sich eine Landschaft mit mehreren lokalen Minima und einem globalen Minimum
erzeugen --- \glqq an ideal playground for testing optimization algorithms\grqq.

\medskip
\textbf{[Ergänzung]} Intuition: Jeder Summand ist ein nach unten gezogener
\glqq Trichter\grqq{} um die Stelle $a_i$; je kleiner $r_i$, desto schmaler und (bei festem
$c_i$) tiefer ist der Trichter, denn am Lochzentrum $x=a_i$ gilt $f \approx -c_i/r_i$
für diesen Term.
\end{defbox}

\begin{bspbox}{Konkrete Shekel-Instanz mit drei Foxholes [S.~48]}
Im Skript wird folgende konkrete Instanz mit $m=3$ Löchern verwendet:
\[
  f(x) \;=\; -\frac{1}{(x-0)^2 + 0.7} \;-\; \frac{1.7}{(x-4)^2 + 0.2} \;-\; \frac{1.1}{(x-7)^2 + 0.4}
\]
Das \textbf{globale Minimum} liegt bei $x = 4$ (tiefstes Loch: gro{\ss}es $c_2 = 1.7$,
kleines $r_2 = 0.2$); die beiden anderen Minima bei $x = 0$ und $x = 7$ sind nur
\emph{lokale} Minima. Wichtig: Diese Information steht uns beim Lauf eines
Optimierungsalgorithmus natürlich \emph{nicht} zur Verfügung --- wir sehen nur die
Funktionswerte und Gradienten an den ausgewerteten Punkten. Für das Studium des
Verhaltens von Optimierungsmethoden ist das Wissen aber nützlich.

\medskip
\textbf{Sensitivität gegenüber dem Startpunkt $x_0$:} Startet Newtons Methode nahe
$x = 0$ oder $x = 7$, konvergiert sie zu einem \emph{lokalen} Minimum --- nicht zum
globalen bei $x = 4$ (vgl.\ Figure~24 [S.~48]: Funktion mit Tälern bei $x = 0, 4, 7$;
die Ableitung kreuzt an jedem Minimum die Null).
\end{bspbox}

\begin{satzbox}{Optimierung als Nullstellensuche: Newton-Update, Gl.~(36) [S.~49]}
Um Newtons Methode zur \emph{Optimierung} (statt Nullstellensuche) einzusetzen, lösen
wir das Nullstellenproblem auf der \emph{Ableitung}: $f'(x) = 0$. Das ist die zentrale
Reformulierung des Optimierungsproblems. Das Newton-Update wird dann
(Original-Gl.~(36)):
\[
  x_{n+1} \;=\; x_n - \frac{f'(x_n)}{f''(x_n)} \tag{36}
\]
\textbf{In eigenen Worten:} Beim gewöhnlichen Newton-Verfahren für Nullstellen von $g$
lautet das Update $x_{n+1} = x_n - g(x_n)/g'(x_n)$. Setzt man $g = f'$, wird daraus
genau Gl.~(36): Zähler $f'$, Nenner $f''$. Randnotiz [S.~48]: Die als gestrichelte Linie
dargestellte Ableitung \emph{ist} die Antwort darauf, wie man Optimierung als
Nullstellensuche formuliert --- \glqq The minima of $f(x)$ correspond to the roots of
$f'(x)$.\grqq{} An Extrema ist $f'(x)$ null; gefunden werden dabei allerdings Minima
\emph{oder} Maxima (Klassifikation: siehe Second-Derivative-Test weiter unten).
\end{satzbox}

\begin{bspbox}{Newton auf der Shekel-Funktion ab $x_0 = 7.2$ (vollständige Rechnung) [S.~49]}
Wir wenden Gl.~(36) auf die obige Drei-Löcher-Instanz an. Die Ableitungen vereinfachen
sich zu (Quotientenregel auf jeden Summanden; das Minuszeichen der Funktion dreht das
Vorzeichen):
\[
  f'(x) = \frac{2x}{(x^2 + 0.7)^2} + \frac{3.4\,(x-4)}{\bigl((x-4)^2 + 0.2\bigr)^2}
        + \frac{2.2\,(x-7)}{\bigl((x-7)^2 + 0.4\bigr)^2}
\]
Auswertung bei $x_0 = 7.2$:
\[
  f'(7.2) \;\approx\; 0.005 + 0.100 + 2.27 \;=\; 2.38
\]
\emph{Anmerkung zur Rundung [S.~49]:} Die drei Summanden ergeben arithmetisch $2.375$;
das PDF rundet zulässig auf $2.38$ --- keine Korrektur erforderlich.
\[
  f''(x) = \frac{2\,(0.7 - 3x^2)}{(x^2 + 0.7)^3}
         + \frac{3.4\,\bigl(0.2 - 3(x-4)^2\bigr)}{\bigl((x-4)^2 + 0.2\bigr)^3}
         + \frac{2.2\,\bigl(0.4 - 3(x-7)^2\bigr)}{\bigl((x-7)^2 + 0.4\bigr)^3}
\]
\[
  f''(7.2) \;\approx\; -0.002 - 0.091 + 7.23 \;=\; 7.14
\]
(Vgl.\ Figure~25 [S.~49]: $f'$ schwarz, $f''$ grau gestrichelt auf der
Shekel-Instanz; \glqq the second derivative is positive at minima, confirming local
curvature\grqq{} --- $f'' > 0$ an den Minima bestätigt die lokale Krümmung.)
Damit lauten die \textbf{Newton-Schritte}:
\[
  x_1 = 7.2 - \frac{2.38}{7.14} \approx 6.87, \qquad
  x_2 \approx 7.02, \qquad
  x_3 \approx 7.00
\]
[Anm.: Exakt nachgerechnet ist $x_3 \approx 6.9907$ ($\approx 6.99$, PDF: $7.00$); das lokale Minimum liegt durch die Überlappung der drei Foxhole-Terme knapp unterhalb von $7$ bei $x \approx 6.9907$. Die Kernaussage (Konvergenz zum lokalen Minimum $x^\circ \approx 7$) bleibt unverändert richtig.]
\textbf{Kernbeobachtung [S.~49]:} Nullstellenverfahren lassen sich zur Optimierung
nutzen, indem man sie auf die \emph{Ableitung} anwendet: $f'(x)$ ist an Extrema null,
dort finden wir also Minima oder Maxima.

\medskip
\textbf{Das fundamentale Problem [S.~50]:} Hier hatten wir Pech --- die Methode wird vom
nahegelegenen Foxhole bei $x^{\circ} \approx 7$ angezogen (zur Notation $x^{\circ}$
siehe unten). Lokale Optimierung garantiert nur das Finden \emph{naher} Extrema.
Selbst mit Konvergenzgarantien: Von $x_0 = 3$ aus hätten wir Glück gehabt, von
$x_0 = -2$ oder $x_0 = 8$ aus verfehlen wir das globale Minimum. --- Genug Motivation
für die globalen Optimierungsstrategien.

\medskip
\emph{Randnotiz Notation [S.~49]:} Der Stern ($*$) ist konventionell für \emph{globale}
Optima reserviert. Zur Unterscheidung lokaler von globalen Minima bieten sich
alternative Notationen wie $x^{\circ}$ (x-circle) oder $x_i$ für das $i$-te lokale
Minimum an.
\end{bspbox}

\subsection{Definitionen: Globale Optimierungsstrategien [S.~51--53]}

\begin{defbox}{Lokale vs.\ globale Optimierung, Gl.~(37) und (38) [S.~51]}
\textbf{Lokale Optimierung (local optimization)} findet ein \emph{nahes} Minimum:
\[
  \mathbf{x}^{*} \;=\; \arg\min_{\mathbf{x} \in \mathcal{N}(\mathbf{x}_0)} f(\mathbf{x}) \tag{37}
\]
wobei $\mathcal{N}(\mathbf{x}_0)$ eine \emph{Nachbarschaft} (neighborhood) des
Startpunkts ist. Randnotiz [S.~51]: $\mathcal{N}$ ist eine \emph{Menge} von Punkten um
$\mathbf{x}_0$, z.\,B.\ $\mathcal{N}(\mathbf{x}_0) = \{\mathbf{x} :
\|\mathbf{x} - \mathbf{x}_0\| < \epsilon\}$ für einen Radius $\epsilon$.

\medskip
\textbf{Globale Optimierung (global optimization)} findet das \emph{beste} Minimum:
\[
  \mathbf{x}^{*} \;=\; \arg\min_{\mathbf{x} \in \Omega} f(\mathbf{x}) \tag{38}
\]
wobei $\Omega$ die gesamte Suchdomäne ist. Das Skript merkt an: \glqq It is easy to see
how the search domain can be inifinitely [sic] large.\grqq{} --- die Suchdomäne kann
unendlich gro{\ss} sein, was die Aufgabe fundamental schwerer macht.
\end{defbox}

\begin{defbox}{Konvexität als Trennlinie (convexity is the dividing line) [S.~51]}
Eine Funktion $f$ hei{\ss}t \textbf{konvex} (convex), wenn
\[
  f\bigl(\lambda x + (1-\lambda)\,y\bigr) \;\le\; \lambda f(x) + (1-\lambda)\,f(y)
  \qquad \text{für alle } \lambda \in [0,1].
\]
\textbf{[Ergänzung]} Anschaulich: Die Verbindungsstrecke (Sehne) zwischen zwei beliebigen
Kurvenpunkten liegt nie unterhalb des Graphen --- die Funktion hat eine einzige
\glqq Schüssel\grqq-Form ohne Nebentäler.

Für konvexe Funktionen ist \emph{jedes lokale Minimum das globale Minimum} ---
lokale Methoden genügen; jede lokale Methode findet das globale Optimum. Die
Schwierigkeit entsteht erst bei \textbf{nicht-konvexer} Optimierung, deren Landschaft
mehrere lokale Minima, Sattelpunkte (saddle points) und Plateaus enthält ---
wie Shekel's Foxholes oder die Loss-Flächen neuronaler Netze.
\end{defbox}

\begin{defbox}{Anziehungsgebiet (basin of attraction) [S.~51, 55]}
Das \textbf{Anziehungsgebiet} eines (lokalen) Minimums ist die Menge aller Startpunkte,
von denen aus ein lokaler Optimierer zu genau diesem Minimum konvergiert.
\textbf{[Ergänzung]} Die Suchdomäne zerfällt dadurch in \glqq Einzugsgebiete\grqq:
Wo immer man startet, rutscht der lokale Optimierer in das Minimum des jeweiligen
Basins. Der Anteil $p$, den das Basin des \emph{globalen} Minimums an der Suchdomäne
hat, bestimmt direkt die Erfolgswahrscheinlichkeit von Random Restarts (Gl.~(39)).
Randnotiz-Übung [S.~55]: Shekel ist insgesamt nicht-konvex, die einzelnen Basins sind
es jedoch; man soll auch eine Funktion notieren, bei der die Basins \emph{nicht} so
leicht zu kartieren sind.
\end{defbox}

\begin{defbox}{Random Restarts [S.~51]}
\textbf{Random Restarts} (zufällige Neustarts) ist die einfachste globale Strategie:
Man startet den lokalen Optimierer wiederholt von zufällig gezogenen Startpunkten der
Suchdomäne und behält die beste gefundene Lösung. Das Skript kommentiert:
\glqq in its notion it feels like brute-forcing again\grqq{} --- vom Ansatz her fühlt es
sich wieder wie Brute-Forcing an. Ablauf siehe Algorithmus~\ref{alg:randomrestarts}
(Original: Algorithmus~3). Randnotiz [S.~51]: Das Training mehrerer neuronaler Netze
mit verschiedenen Seeds ist \emph{implizites} Random Restarts --- für gro{\ss}e Modelle
jedoch nicht praktikabel.
\end{defbox}

\begin{algorithm}[H]
\caption{Random Restarts (Original: Algorithmus~3 [S.~51])}
\label{alg:randomrestarts}
\begin{algorithmic}[1]
\Require Suchdomäne $\Omega$, lokaler Optimierer \textsc{LocalOptimize}, Anzahl Restarts $K$
\State $x_{\text{best}} \gets \textsc{None}$;\quad $f_{\text{best}} \gets +\infty$
\For{$k = 1$ \textbf{to} $K$}
  \State $x_0 \gets \textsc{RandomPoint}(\Omega)$
  \State $x^{*} \gets \textsc{LocalOptimize}(x_0)$
  \If{$f(x^{*}) < f_{\text{best}}$}
    \State $x_{\text{best}} \gets x^{*}$;\quad $f_{\text{best}} \gets f(x^{*})$
  \EndIf
\EndFor
\State \Return $x_{\text{best}}$
\end{algorithmic}
\end{algorithm}

\begin{satzbox}{Erfolgswahrscheinlichkeit von Random Restarts, Gl.~(39) [S.~51]}
Hat das Anziehungsgebiet des globalen Minimums die Wahrscheinlichkeit $p$ (Anteil an
der Suchdomäne, aus der ein zufälliger Startpunkt ins globale Basin fällt), dann gilt
mit $K$ Restarts:
\[
  P(\text{find global}) \;=\; 1 - (1 - p)^{K} \tag{39}
\]
\textbf{In eigenen Worten:} Ein einzelner Restart \emph{verfehlt} das globale Basin mit
Wahrscheinlichkeit $1-p$. Da die Restarts unabhängig sind, verfehlen alle $K$ Versuche
gemeinsam mit Wahrscheinlichkeit $(1-p)^K$; das Komplement davon ist die
Erfolgswahrscheinlichkeit. Zahlenbeispiel des Skripts: Für $p = 0.1$ und $K = 20$ ist
$P = 1 - 0.9^{20} \approx 0.88$.

\medskip
\textbf{Auflösen nach $K$} (Randnotiz-Hint [S.~54]): Aus $P = 1-(1-p)^K$ folgt
\[
  K \;=\; \frac{\ln(1-P)}{\ln(1-p)}.
\]
\textbf{Praxisproblem:} Üblicherweise ist $p$ vorab \emph{nicht} bekannt. Heuristik für
$K$ ohne Kenntnis von $p$: $K$ erhöhen, bis sich die beste Lösung stabilisiert (keine
Verbesserung nach mehreren Restarts) --- siehe Patience.
\end{satzbox}

\begin{defbox}{Patience (Geduld-Hyperparameter) [S.~51]}
\textbf{Patience} ist der übliche Hyperparameter der $K$-Heuristik: Er kontrolliert,
wie viele Restarts \emph{ohne Verbesserung} abgewartet werden, bevor gestoppt wird.
\textbf{[Ergänzung]} Dieselbe Idee kennt man aus dem Deep Learning als
Early-Stopping-Patience: Statt eine feste Versuchszahl vorab festzulegen, beendet man
adaptiv, wenn weitere Versuche längere Zeit nichts mehr verbessern.
\end{defbox}

\begin{defbox}{Basin Hopping [S.~52]}
\textbf{Basin Hopping} erkundet die Landschaft durch \emph{Abwechseln} von lokaler
Optimierung und zufälligen Sprüngen (random jumps). Das ist verschieden von Random
Restarts, das die Suche jeweils \emph{komplett zurücksetzt}: Basin Hopping erlaubt das
Entkommen aus lokalen Minima, während weiterhin lokale Optimierung genutzt wird, um
bessere Lösungen zu finden --- man springt vom gerade gefundenen lokalen Minimum aus
weiter, statt blind neu zu starten. Die \textbf{Sprungverteilung} (jump distribution)
wird als $\delta \sim \mathcal{D}$ spezifiziert --- z.\,B.\ eine bei null zentrierte
Gau{\ss}-Verteilung oder eine Gleichverteilung über einen bestimmten Bereich. Ablauf:
Algorithmus~\ref{alg:basinhopping} (Original: Algorithmus~4).
\end{defbox}

\begin{algorithm}[H]
\caption{Basin Hopping (Original: Algorithmus~4 [S.~52])}
\label{alg:basinhopping}
\begin{algorithmic}[1]
\Require Startpunkt $x_0$, lokaler Optimierer \textsc{LocalOptimize}, Sprungverteilung für $\delta$, Iterationen $K$
\State $x_{\text{best}} \gets x_0$;\quad $f_{\text{best}} \gets f(x_0)$
\State $x_{\text{current}} \gets x_0$
\For{$k = 1$ \textbf{to} $K$}
  \State $x^{*} \gets \textsc{LocalOptimize}(x_{\text{current}})$
  \If{$f(x^{*}) < f_{\text{best}}$}
    \State $x_{\text{best}} \gets x^{*}$;\quad $f_{\text{best}} \gets f(x^{*})$
  \EndIf
  \State Sprung ziehen: $\delta \sim \mathcal{D}$
  \State $x_{\text{current}} \gets x^{*} + \delta$ \Comment{Random jump}
\EndFor
\State \Return $x_{\text{best}}$
\end{algorithmic}
\end{algorithm}

\begin{defbox}{Simulated Annealing (simuliertes Abkühlen) [S.~52]}
\textbf{Simulated Annealing} kombiniert Basin Hopping mit einem
\emph{probabilistischen Akzeptanzkriterium} für Kandidatenpunkte: Zu Beginn wird ein
Sprung in eine \emph{schlechtere} Lösung mit Wahrscheinlichkeit
\[
  \exp\!\left(-\frac{\Delta f}{T}\right)
\]
akzeptiert, wobei $T$ (\emph{Temperatur}) über die Zeit abnimmt. Randnotiz [S.~52]:
$\Delta f$ ist der \emph{Anstieg} des Zielfunktionswerts beim Übergang von der
aktuellen Lösung zu einer schlechteren Kandidatenlösung. Später wird der Algorithmus
konservativer und akzeptiert nur noch Kandidaten, die das Ziel verbessern --- so kann
er zu einem Optimum konvergieren. \textbf{[Ergänzung]} Intuition: Bei hohem $T$ ist
$\Delta f / T$ klein, die Akzeptanzwahrscheinlichkeit also nahe~1 --- viel Exploration;
bei $T \to 0$ geht sie für jede Verschlechterung gegen~0 --- reine Ausbeutung
(Exploitation). Der Name stammt vom kontrollierten Abkühlen (Annealing) in der
Metallurgie.
\end{defbox}

\begin{defbox}{Stochastische Methoden (stochastic methods) und Noisy Newton [S.~52--53]}
\textbf{Stochastische Methoden} fügen Rauschen hinzu, indem die Loss-Landschaft
\emph{approximiert} wird, um lokalen Optima zu entkommen. Beispielaufgabe des Skripts:
die 1D-Shekel-Funktion als Funktion zweiter Ordnung approximieren, indem man drei
Punkte auf der Foxhole-Kurve wählt (vgl.\ Figure~26 [S.~53]: die Parabel durch drei
Stützpunkte erfasst den allgemeinen Trend, glättet aber die einzelnen Foxholes weg).

\medskip
\textbf{Noisy Newton} (Randnotiz [S.~52]): folgt derselben Hopping-Idee, indem
Gau{\ss}-Rauschen direkt zum Newton-Schritt addiert wird:
$\mathbf{x}_{n+1} = \ldots + \boldsymbol{\xi}_n$.

\medskip
\textbf{The take away [S.~52--53]:} Die Loss-Landschaft ist \emph{nicht statisch},
sondern kann konstruiert und geformt werden (loss landscape engineering). Durch
Approximation der Loss-Landschaft auf Mini-Batches (jeweils andere Punktauswahl)
erhält man eine \emph{verrauschte Schätzung} der wahren Loss-Landschaft, die hilft,
lokalen Minima zu entkommen. Die Loss-Landschaft wird natürlich auch stark durch die
Wahl der Loss-Funktion selbst beeinflusst. Aber: Der Newton-Schritt ist sehr
\emph{rauschempfindlich} --- später werden Ansätze folgen, die damit besser umgehen.
\end{defbox}

\subsection{Newtons Methode zum Finden von Minima [S.~53]}

\begin{satzbox}{Newton-Update für Minima und Second-Derivative-Test, Gl.~(40) [S.~53]}
Die im Hands-on gesehene Reformulierung von Optimierung als Nullstellensuche erlaubt
Newtons Methode zum Finden von Optima (Original-Gl.~(40), identisch zu Gl.~(36)):
\[
  x_{n+1} \;=\; x_n - \frac{f'(x_n)}{f''(x_n)} \tag{40}
\]
\textbf{Minima, Maxima und Sattelpunkte:} Newtons Methode (auf $f'$ angewandt) findet
Punkte mit $f'(x) = 0$ \emph{[Korrektur ggü.\ Original, S.~53: das PDF druckt
\glqq Newton's method finds points where $f''(x) = 0$\grqq; sachlich findet Newton auf
$f'$ angewandt Punkte mit $f'(x) = 0$ (Extremstellen); erst der anschlie{\ss}ende
Second-Derivative-Test klassifiziert sie über $f''$]} --- diese Punkte können aber
Minima, Maxima oder Sattelpunkte sein. Der \textbf{Test der zweiten Ableitung}
(second derivative test) unterscheidet:
\begin{enumerate}[label=(\alph*),itemsep=1pt]
  \item $f''(x^{*}) > 0$: lokales \textbf{Minimum} (\glqq curve bends upward\grqq{} ---
        die Kurve biegt nach oben);
  \item $f''(x^{*}) < 0$: lokales \textbf{Maximum} (\glqq curve bends downward\grqq{} ---
        die Kurve biegt nach unten);
  \item $f''(x^{*}) = 0$: \textbf{nicht schlüssig} (inconclusive) --- möglicher
        Sattelpunkt oder Wendepunkt (saddle point / inflection point).
\end{enumerate}
\end{satzbox}

\begin{satzbox}{Newton Richtung Minima lenken: der $|f''|$-Trick, Gl.~(41) [S.~53--54]}
Modifiziertes Update mit dem \emph{Betrag} der zweiten Ableitung im Nenner
(Original-Gl.~(41)):
\[
  x_{n+1} \;=\; x_n - \frac{f'(x_n)}{\bigl|f''(x_n)\bigr|} \tag{41}
\]
So wird sichergestellt, dass wir uns \emph{immer} in Richtung abnehmenden $f(x)$
bewegen: Der Schritt geht stets in Richtung des negativen Gradienten
(\glqq downhill so to speak\grqq).
\textbf{In eigenen Worten / [Ergänzung]:} Beim reinen Newton-Update (Gl.~(40))
bestimmt das \emph{Vorzeichen} von $f''$ die Schrittrichtung mit. In einer Region mit
$f'' < 0$ (Kurve biegt nach unten, also nahe einem Maximum) zeigt der Newton-Schritt
\emph{bergauf} --- Newton wird vom Maximum angezogen. Nimmt man $|f''|$, bleibt die
adaptive Schrittweite $1/|f''|$ erhalten, aber die Richtung ist immer
$-\operatorname{sign}(f'(x_n))$, also bergab. Geometrisch (Figures~27/28 [S.~54],
Beispiel am Punkt $x = 6.1$): Die Tangente an $f'$ hat dort Steigung $f''(6.1) \approx -2.98$
[Korrektur ggü.\ Original, S.~54: Die Captions von Fig.~27/28 drucken $\approx \mp 2.66$ ---
das ist $f''$ bei $x \approx 6.0$ ($f''(6.0) \approx -2.63$); bei $x = 6.1$ ist
$f'' \approx -2.98$]; der Ersatz $f'' \to |f''|$ klappt die Steigung auf $\approx +2.98$ um --- das
Vorzeichen der Krümmung wird entfernt, Abwärtsschritte werden erzwungen und die Gefahr
sinkt, auf ein lokales Maximum zuzulaufen. Randnotiz-Denkfrage [S.~53]: Wie müsste das
Update modifiziert werden, um stattdessen \emph{Maxima} zu finden?

\medskip
\textbf{Randnotiz \glqq In deep learning\grqq{} [S.~53] --- die vier DL-Hoffnungen:}
Im Deep Learning verwendet man typischerweise lokale Methoden (SGD, Adam), aber
\emph{hofft}, dass: (1) die meisten lokalen Minima ungefähr gleich gut sind,
(2) Überparametrisierung (overparameterization) schlechte Minima selten macht,
(3) Mini-Batch-Rauschen hilft, scharfen Minima zu entkommen,
(4) adaptive Lernraten helfen, Plateaus zu durchqueren.
\end{satzbox}

\subsection{Durchgerechnete Beispiele: Strategien auf konkreten Zahlen [S.~54--55]}

\begin{bspbox}{Random Restarts: Erfolgswahrscheinlichkeit und nötige Restart-Zahl [S.~54--55]}
\textbf{Aufgabe:} Eine Funktion hat 5 lokale Minima; das Anziehungsgebiet des globalen
Minimums deckt 15\,\% der Suchdomäne ab ($p = 0.15$). Wie gro{\ss} ist die
Wahrscheinlichkeit, das globale Minimum mit $K = 10$ Random Restarts zu finden?

\medskip
\textbf{Rechnung mit Gl.~(39):}
\[
  P \;=\; 1 - (1 - 0.15)^{10} \;=\; 1 - 0.85^{10} \;=\; 1 - 0.1969 \;=\; 0.803
\]
Etwa 80\,\% Chance --- \glqq not too bad, but far from certain\grqq{} (nicht schlecht,
aber weit von Gewissheit entfernt).

\medskip
\textbf{Folgefrage: Wie viele Restarts brauchen wir für 99\,\%?} Mit der nach $K$
aufgelösten Formel (Randnotiz-Hint [S.~54]: $K = \ln(1-P)/\ln(1-p)$):
\[
  K \;=\; \frac{\ln(1 - 0.99)}{\ln(1 - 0.15)} \;=\; \frac{\ln(0.01)}{\ln(0.85)}
    \;=\; \frac{-4.605}{-0.1625} \;\approx\; 28.3
\]
[S.~55] Da $K$ ganzzahlig sein muss und das Ziel erreicht werden soll, wird
aufgerundet: \textbf{$K = 29$ Restarts.}

\medskip
\textbf{[Ergänzung]} Beachte die ungünstige Skalierung: Je kleiner das globale Basin
($p \downarrow$), desto stärker wächst $K$, denn $\ln(1-p) \approx -p$ für kleine $p$,
also $K \approx -\ln(1-P)/p$ --- der Aufwand wächst etwa wie $1/p$.
(Anschlussübung des Skripts [S.~55]: $K$ für $p = 0.05$ und $p = 0.01$ beim selben
99-\%-Ziel berechnen --- siehe Übungsteil.)
\end{bspbox}

\subsection{Wissenswertes aus Randnotizen und Fu{\ss}noten}
\begin{itemize}[itemsep=2pt]
  \item \textbf{Loss-Landschaften neuronaler Netze [S.~47]:} Li et al.\ (2018,
        Fu{\ss}note~8: \glqq Visualizing the loss landscape of neural nets\grqq,
        NeurIPS) zeigen das komplexe Terrain der Optimierung neuronaler Netze ---
        verschiedene Initialisierungen führen zu verschiedenen Endlösungen.
  \item \textbf{Fu{\ss}note~9 [S.~48]:} J.~Shekel, \glqq Test functions for multimodal
        search techniques\grqq, 1971 --- Ursprung der Shekel-Funktion.
  \item \textbf{Implizite Random Restarts [S.~51]:} Das Training mehrerer Netze mit
        verschiedenen Seeds entspricht implizitem Random Restarts; für gro{\ss}e
        Modelle nicht praktikabel.
  \item \textbf{Code [S.~54]:}
        \url{https://github.com/Quillstacks/lecturecode_numericalmethods.git}
  \item \textbf{Teaser [S.~56]:} Gradientenabstieg neigt zum Steckenbleiben ---
        \glqq so far our answer was merely better than brute-forcing\grqq. Nächstes
        Kapitel: die Loss-Landschaft \emph{glätten}, um die Robustheit des
        Gradientenabstiegs zu erhöhen. Randnotiz-Denkfrage: Was, wenn wir einer
        \emph{Hochfrequenz-Version} von Shekel's Foxholes gegenüberstehen --- welches
        Problem entsteht? (Antwort in Thema~6.)
\end{itemize}

\subsection{Übungsaufgaben und Selbstreflexionsfragen des Skripts}

\textbf{Übungsaufgaben [S.~48--55] (nur Fragen):}
\begin{itemize}[itemsep=2pt]
  \item Newton auf der 1D-Shekel-Funktion von einem \emph{anderen} Startpunkt als
        $x_0 = 7.2$ anwenden und beobachten, wohin die Methode konvergiert [S.~48].
  \item Die 1D-Shekel-Funktion als Funktion zweiter Ordnung approximieren, indem drei
        Punkte auf der Foxhole-Kurve gewählt werden [S.~52].
  \item Berechne $K$ für $p = 0.05$ und $p = 0.01$ beim selben 99-\%-Ziel [S.~55].
  \item \emph{On your machine} [S.~55]: Eigene 1D-Shekel-Funktion designen (anfangs
        einfach halten, später komplexer machen); die bisher gelernten lokalen
        Optimierungsmethoden darauf anwenden; zeigen, wo Newtons Methode (wieder)
        scheitert; verschiedene $x_0$ erkunden und das Konvergenzverhalten beobachten;
        über Loss-Landschaft und Basin-Grenzen nachdenken.
  \item \emph{Directing the search towards minima} [S.~55]: Der $|f''|$-Trick dreht
        das Vorzeichen der Krümmung im Nenner um und erzwingt Bergab-Schritte ---
        was bedeutet das geometrisch für die Tangentenkonstruktion? Danach im Code
        ausprobieren.
  \item \emph{Escaping Local Minima} [S.~55]: Random Restarts und Basin Hopping auf
        der eigenen 1D-Shekel-Funktion einsetzen; vergleichen, wie viele
        Funktionsauswertungen jede Methode zum Finden des globalen Minimums braucht.
        Wie anfällig sind sie für das Steckenbleiben in lokalen Minima? Wie
        beeinflusst die Wahl der Schrittweitenverteilung die Performance von Basin
        Hopping?
  \item \emph{Noise and landscape engineering} [S.~55]: Die Loss-Landschaft ist nicht
        statisch, sondern zu konstruieren und zu formen; durch Approximation auf
        Mini-Batches (verschiedene Punktauswahlen) \ldots\ [Anm.: Der Satz endet im
        Original abrupt mit Komma --- unvollständig im PDF, Inhalt nur teilweise
        rekonstruierbar.]
  \item \emph{Code exercise} [S.~55]: Random Restarts und Basin Hopping auf der
        1D-Shekel-Funktion implementieren; vergleichen, wie viele
        Funktionsauswertungen jede Methode braucht, um das globale Minimum bei
        $x \approx 4$ zu finden.
  \item \emph{Basin of Attraction} (Randnotiz [S.~55]): Die Plots der Shekel-Funktion
        erneut ansehen, die Anziehungsgebiete jedes lokalen Minimums markieren und
        ihre relativen Grö{\ss}en schätzen. Welches gehört zum globalen Minimum? Wie
        hängt das mit der Erfolgswahrscheinlichkeit von Random Restarts zusammen?
  \item \emph{Non-Convex Basins} (Randnotiz [S.~55]): Shekel ist insgesamt
        nicht-konvex, die Basins jedoch schon --- eine Funktion notieren, bei der die
        Anziehungsgebiete nicht so leicht zu kartieren sind.
  \item \emph{Schrittweiten-Heuristik} (Randnotiz [S.~55]): Gibt es einen klugen Weg,
        die Schrittweitenverteilung während der Optimierung automatisch anzupassen?
        Was wäre eine gute Heuristik? Wie verhindert man katastrophale Sprünge?
        Implementieren.
  \item \emph{Basins of Attraction kartieren} [S.~55]: Für die 1D-Shekel-Funktion
        (Annahme: ein lokaler Optimierer konvergiert immer zum nächstgelegenen
        Foxhole [Anm.: im PDF folgt ein Fu{\ss}notenrest \glqq foxhole.2\grqq]):
        \begin{enumerate}[label=(\alph*),itemsep=1pt]
          \item Die Anziehungsgebiete der drei Foxholes bei $x = 0, 4, 7$ grob
                skizzieren --- wo liegen die Basin-Grenzen?
          \item Wenn Basin Hopping einen Sprung $\delta \sim \mathrm{Uniform}(-3, 3)$
                vom lokalen Minimum $x^{\circ} = 7$ aus nutzt: Mit welcher
                Wahrscheinlichkeit landet ein einzelner Sprung im Basin des globalen
                Minimums?
          \item Warum könnte Basin Hopping das globale Minimum hier schneller finden
                als Random Restarts?
        \end{enumerate}
\end{itemize}

\textbf{Selbstreflexionsfragen [S.~56]:}
\begin{itemize}[itemsep=2pt]
  \item Warum scheitert lokale Optimierung auf nicht-konvexen Landschaften wie
        Shekel's Foxholes?
  \item Wie skaliert der nötige Aufwand bei Random Restarts mit der Grö{\ss}e des
        globalen Basins?
  \item Wie unterscheidet sich Basin Hopping von Random Restarts? [Originalfrage
        grammatisch fehlerhaft: \glqq How do basin hopping differently from random
        restarts?\grqq{} [sic]]
  \item Was sagt uns $f''(x_n)$ über die Loss-Landschaft, was charakterisiert es, und
        wie können wir es nutzen?
  \item Wie hilft die Approximation der Loss-Landschaft mit verrauschten Schätzungen
        (Mini-Batches) beim Entkommen aus lokalen Optima?
  \item Wähle $x_0$ so, dass es im Basin eines lokalen Minimums liegt; finde
        verschiedene Wege, mit denen die Optimierungsmethode entkommen kann --- wie
        effektiv ist was?
\end{itemize}

\subsection{Recap of Key Concepts [S.~56]}
\begin{itemize}[itemsep=2pt]
  \item Lokale Optimierungsmethoden (wie Newtons Methode) können auf nicht-konvexen
        Funktionen in lokalen Minima stecken bleiben.
  \item Globale Optimierungsstrategien (Random Restarts, Basin Hopping, Simulated
        Annealing, Stochastizität) erkunden den Suchraum gezielt breiter, um das
        globale Optimum zu finden.
  \item Die Modifikation von Optimierungsmethoden mit Krümmungsinformation $f''(x)$
        erlaubt es, die Abstiegsrichtung genauer zu spezifizieren.
  \item Die Loss-Landschaft \glqq gehört uns\grqq{} --- sie ist durch die Wahl der
        Loss-Funktion und durch Rauschen (etwa stochastische Gradienten) zu
        konstruieren (loss landscape engineering).
\end{itemize}

\begin{merkbox}
\textbf{Typische Fehlerquellen:}
\begin{itemize}[itemsep=1pt]
  \item \textbf{Newton auf $f$ statt $f'$:} Für Optimierung wird das
        Nullstellenproblem auf der \emph{Ableitung} gelöst --- Update
        $x_{n+1} = x_n - f'(x_n)/f''(x_n)$ (Gl.~(36)/(40)), nicht $f/f'$.
  \item \textbf{$f'(x^{*})=0$ hei{\ss}t nicht automatisch Minimum:} Newton auf $f'$
        findet Extremstellen ($f'=0$), die auch Maxima oder Sattelpunkte sein können
        --- immer mit dem Second-Derivative-Test klassifizieren. (Achtung: Das
        Original druckt auf S.~53 fälschlich \glqq $f''(x)=0$\grqq{} --- korrekt ist
        $f'(x)=0$.)
  \item \textbf{Vorzeichen von $f''$ ignoriert:} In Regionen mit $f''<0$ läuft das
        reine Newton-Update bergauf zum Maximum; der $|f''|$-Trick (Gl.~(41))
        erzwingt Bergab-Schritte.
  \item \textbf{Wahrscheinlichkeiten falsch kombiniert:} Erfolgswahrscheinlichkeit ist
        $1-(1-p)^K$ --- \emph{nicht} $K \cdot p$ (das kann $>1$ werden und ignoriert
        die Unabhängigkeit der Versuche).
  \item \textbf{$K$ abrunden:} Beim Auflösen nach $K$ stets \emph{aufrunden}
        ($28.3 \to 29$), sonst wird das Wahrscheinlichkeitsziel verfehlt.
  \item \textbf{Random Restarts und Basin Hopping verwechseln:} Restarts setzen die
        Suche komplett zurück; Basin Hopping springt vom zuletzt gefundenen lokalen
        Minimum aus weiter ($x_{\text{current}} \gets x^{*} + \delta$).
\end{itemize}
\textbf{Merksätze:}
\begin{itemize}[itemsep=1pt]
  \item \glqq Die Minima von $f$ sind die Nullstellen von $f'$.\grqq
  \item Konvexität ist die Trennlinie: konvex $\Rightarrow$ lokal = global; die ganze
        Mühe dieses Kapitels gilt dem nicht-konvexen Fall.
  \item Erst finden ($f'=0$), dann klassifizieren ($f''$): positiv = Tal, negativ =
        Gipfel, null = unklar.
  \item Simulated Annealing: hei{\ss} = experimentierfreudig
        ($\exp(-\Delta f/T) \approx 1$), kalt = konservativ.
  \item Rauschen ist ein Werkzeug, kein Feind: Mini-Batch-Rauschen hilft beim
        Entkommen --- aber der Newton-Schritt verträgt Rauschen schlecht.
\end{itemize}
\end{merkbox}
